\noindent302\hfill\textbf{NACHLASS.}
\begin{center}[4.]\\[4pt]
[Normalgleichungen zu den Winkelgleichungen.]\end{center}

[Bezeichnen $(1.2)$, $(1.3)$, $(2.3)$ usw. die Verbesserungen der Richtungsbeobachtungen
$1.2$, $1.3$, $2.3$ usw., so ergeben sich für die 51 Dreiecke
die folgenden Bedingungsgleichungen:
\[
\begin{aligned}
+(1.2)-(2.1)+(2.3)-(3.2)+(3.1)-(1.3)-1{,}436&=0\\
+(2.4)-(4.2)+(4.3)-(3.4)+(3.2)-(2.3)+1{,}294&=0
\end{aligned}
\]
\begin{center}usw.\end{center}

Von diesen sind jedoch nur 43 voneinander unabhängig. Die Korrelaten
der Bedingungsgleichungen seien entsprechend den Dreiecksnummern $(1)$, $(2)$,
$(3)$ usw. Haben die Richtungen gleiche Gewichte und bleiben die
Bedingungsgleichungen aus den Seitenverhältnissen unberücksichtigt,
so gilt:
\[
\begin{aligned}
(1.2)&=-(2.1)=+(1)\\
(1.3)&=-(3.1)=-(1)\\
(2.3)&=-(3.2)=+(1)-(2)\\
(2.4)&=-(4.2)=+(2)\\
(3.4)&=-(4.3)=-(2)-(3)
\end{aligned}
\]
\begin{center}usw.\end{center}

Nur die Richtung $24.26$ (Bottel---Steinberg) hat das Gewicht $\frac14$ erhalten.
Ihre Verbesserung ist demnach
\[
(24.26)=4(36)-4(44).
\]

Setzt man diese Werte in die Bedingungsgleichungen ein, so erhält man
die Normalgleichungen:
\[
\begin{aligned}
6(1)-2(2)\qquad-1{,}436&=0\\
-2(1)+6(2)+2(3)+1{,}294&=0
\end{aligned}
\]
\begin{center}usw.\end{center}

Da hier nur die Korrelaten zu berechnen sind, können sämtliche
Gleichungen durch 2 geteilt werden. Damit ergeben sich die folgenden
Gleichungen zur Bestimmung der Korrelaten.]

\newpage
% Carl Friedrich Gauss, Werke IX, printed 303–304, Article [4] continuation.
\noindent\textbf{ZUR NETZAUSGLEICHUNG.}\hfill303
\begingroup
\renewcommand{\arraystretch}{1.16}
\begin{longtable}{@{}r@{\quad}l@{}}
1.&$0=-0{,}718+3(1)-(2)$\\
2.&$0=+0{,}647-(1)+3(2)+(3)$\\
3.&$0=+0{,}931+(2)+3(3)-(4)-(5)$\\
4.&$0=-0{,}340-(3)+3(4)$\\
5.&$0=-0{,}331-(3)+3(5)-(6)$\\
6.&$0=-0{,}032-(5)+3(6)-(7)-(8)$\\
(7.&$0=+0{,}233.5-(6)+3(7)+(8)+(9)-(10))$\\
8.&$0=+0{,}510.5-(6)+(7)+3(8)-(9)+(10)$\\
9.&$0=-0{,}442+(7)-(8)+3(9)+(10)-(11)$\\
10.&$0=-0{,}164.5-(7)+(8)+(9)+3(10)-(11)$\\
11.&$0=+0{,}600-(9)-(10)+3(11)-(12)$\\
12.&$0=-0{,}164-(11)+3(12)-(13)-(14)$\\
(13.&$0=-0{,}684-(12)+3(13)+(14)-(15)-(16)-(19))$\\
14.&$0=-0{,}569.5-(12)+(13)+3(14)+(16)+(17)-(18)+(19)$\\
(15.&$0=+0{,}886.5-(13)+3(15)+(16)-(17)+(20))$\\
16.&$0=+0{,}521-(13)+(14)+(15)+3(16)+(17)-(18)-(19)-(20)$\\
17.&$0=-0{,}740.5+(14)-(15)+(16)+3(17)-(18)+(20)-(21)$\\
18.&$0=-0{,}918.5-(14)-(16)-(17)+3(18)-(27)-(28)$\\
19.&$0=-0{,}406.5-(13)+(14)-(16)+3(19)+(20)$\\
20.&$0=-0{,}375+(15)-(16)+(17)+(19)+3(20)-(21)$\\
21.&$0=-0{,}672-(17)-(20)+3(21)-(22)-(23)$\\
22.&$0=+0{,}752.5-(21)+3(22)+(23)-(24)-(25)+(38)$\\
23.&$0=-0{,}228-(21)+(22)+3(23)+(25)+(26)-(29)-(38)$\\
24.&$0=+0{,}660-(22)+3(24)+(25)-(26)+(39)$\\
25.&$0=+0{,}108-(22)+(23)+(24)+3(25)+(26)-(29)+(38)-(39)$\\
26.&$0=-0{,}688.5+(23)-(24)+(25)+3(26)-(29)+(39)$\\
27.&$0=-1{,}577.5-(18)+3(27)+(28)-(29)+(30)+(31)$\\
28.&$0=-0{,}914-(18)+(27)+3(28)-(32)+(33)$\\
29.&$0=-0{,}600.5-(23)-(25)-(26)-(27)+3(29)-(30)-(31)$\\
30.&$0=-0{,}378+(27)-(29)+3(30)+(31)-(32)-(34)+(35)-(40)$\\
31.&$0=-0{,}900+(27)-(29)+(30)+3(31)-(35)-(37)+(40)$\\
32.&$0=+0{,}290-(28)-(30)+3(32)-(33)+(34)-(35)-(41)$\\
33.&$0=+0{,}061.5+(28)-(32)+3(33)+(34)-(36)-(41)$\\
(34.&$0=-0{,}179-(30)+(32)+(33)+3(34)-(35)-(36)+(41)-(42))$\\
\end{longtable}
\endgroup
\newpage
\noindent304\hfill\textbf{NACHLASS.}
\begingroup
\renewcommand{\arraystretch}{1.16}
\begin{longtable}{@{}r@{\quad}l@{}}
(35.&$0=-0{,}270.5+(30)-(31)-(32)-(34)+3(35)+(37)+(40)-(43))$\\
36.&$0=-0{,}533.5-(33)-(34)+\frac92(36)+(37)-\frac52(44)$\\
37.&$0=-0{,}891-(31)+(35)+(36)+3(37)-(45)$\\
(38.&$0=+1{,}088+(22)-(23)+(25)+3(38)-(39))$\\
(39.&$0=-0{,}136+(24)-(25)+(26)-(38)+3(39))$\\
40.&$0=-0{,}792.5-(30)+(31)+(35)+3(40)-(43)$\\
41.&$0=-0{,}530.5-(32)-(33)+(34)+3(41)-(42)$\\
42.&$0=+0{,}684.5-(34)-(41)+3(42)-(43)-(44)$\\
43.&$0=+1{,}421-(35)-(40)-(42)+3(43)+(45)-(46)$\\
(44.&$0=-1{,}819-\frac52(36)-(42)+\frac92(44)+(45))$\\
45.&$0=+0{,}194.5-(37)+(43)+(44)+3(45)-(46)$\\
46.&$0=-1{,}470.5-(43)-(45)+3(46)-(47)$\\
47.&$0=-0{,}137.5-(46)+3(47)-(48)$\\
48.&$0=-0{,}649.5-(47)+3(48)-(49)$\\
49.&$0=-2{,}460-(48)+3(49)-(50)$\\
50.&$0=-0{,}774.5-(49)+3(50)-(51)$\\
51.&$0=+0{,}020-(50)+3(51).$\\
\end{longtable}
\endgroup

[Die eingeklammerten Normalgleichungen entsprechen den 8 abhängigen
Bedingungsgleichungen. Bei der Lösung sind die Korrelaten $(7)$, $(13)$, $(15)$,
$(34)$, $(35)$, $(38)$, $(39)$, $(44)$ gleich null zu setzen.]
